\section{K-035: no odd grading inertia anywhere on the rank-two component}

The normalization computation K-033 and the upstairs bound K-034 did
not determine whether an odd grading subgroup could permute branches
over the conductor.  The following result settles that question for
the actual component of K-033, without proving normality or computing
its conductor.  It does not settle the other components in K-031.

\begin{theorem}[K-035]
Let $Y\subset X_{10}^{\mu_{512}}$ be the four-dimensional reduced
irreducible component constructed in K-033, and put $M=512$.
Every nonzero point of $Y$ has grading stabilizer either $\mu_M$
or $\mu_{2M}$.  In particular,
\[
                     Y\cap X_{10}^{\mu_{3M}}=\{0\}
                       \quad\hbox{on closed points}.    \tag{K-035a}
\]
The same holds for the two cyclic translates of $Y$.

Consequently every finitely generated graded coherent module supported
on these components has a three-integral Laurent expansion at a
primitive $M$th root.  Their entire cyclically transported coherent
contribution is zero modulo three in the principal-part reduction of
K-031.  This includes their conductors; it is not restricted to a
normal or smooth open subset.
\end{theorem}

We prove the theorem by controlling full stabilizers, rather than
deducing it from the normalization Hilbert series.

\subsection*{The fixed model and its exact balance data}

Use the dimension-nine tensor $T_9$ of K-020, with selected edge
$e_0=(8,6,6)$ in zero-based coordinates.  Its support $\Sigma$ has
$25$ independent incidence columns, is free, and has a positive
margin-one mass vector $p$.  The selected mass is $293/512$.
K-020 gives its full stabilizer $S_9\times\mu_{512}$ and its Lie
stabilizer $\operatorname{Lie}S_9$.  These are previously checked
inputs, not inferred merely from the support determinant here.

Let $h$ be the order-$512$ element in K-033, and let
\[
 U=V_{10}^h,
 \qquad L=C_{GL_{10}^3}(h),
 \qquad C=L\cap SL_{10}^3.
\]
In each factor the selected old coordinate and the new coordinate
$9$ form the unique two-dimensional eigenspace.  The other eight
eigenspaces are lines.  The exact spectra, as exponents modulo $512$,
are
\[
\begin{aligned}
 A:&\ (36,79,228,86,132,100,74,98,486,486),\\
 B:&\ (414,457,438,450,286,412,158,416,0,158),\\
 C:&\ (488,19,26,476,510,500,380,126,0,380).
\end{aligned}                                             \tag{K-035b}
\]
The determinant-product exponent is $219$, which is odd.  The first
and third spectra contain no antipodal pair.  The complete fixed
space has exactly the $37$ coordinates described in K-033a.

Every nonzero point of $Y$ has a closed $C$-orbit representative in
$U$: the finite normalization map of K-032, specialized in K-033,
is surjective.  Choose a minimum-norm representative for that closed
orbit.  The moment-map argument in K-032 makes it a full
special-linear critical tensor.  Multiplying it by a positive scalar
does not change its stabilizer, so we henceforth normalize all three
reduced Hermitian matrices to the identity.  Write this tensor as $T$.

Collapse each two-dimensional selected block onto its old coordinate.
The aggregate squared norms in the $25$ old edge groups are forced,
by independence of the old incidence columns, to be
\[
                 p_e\ (e\ne e_0),\qquad p_{e_0}+1\ (e=e_0).
                                                               \tag{K-035c}
\]
Thus every old group is nonzero.  Nineteen groups are scalar; the
remaining six groups are the core cube and five vectors.  Their
norm-squared totals, with common denominator $512$, are
\[
                       (805,55,164,219,200,19),           \tag{K-035d}
\]
where the vector axes are, in order, $A,A,B,C,C$.

\subsection*{Centralizers and a positive-dimensional bridge}

Put
\[
 \widetilde H=\operatorname{Stab}_{GL_{10}^3}(T),\qquad
 S=\{(aI,bI,cI):abc=1\},\qquad H=\widetilde H/S,
 \qquad J=\widetilde H\cap L.
\]
All natural factor representations below are representations of
$\widetilde H$, not of the projective quotient $H$.

First, $\chi(\widetilde H)$ is finite.  A homogeneous special-linear
invariant nonzero at $T$ forces a positive power of $\chi$ to be
one on its stabilizer.  Moreover, $\widetilde H$ is closed under
Hermitian adjoints.  To see this directly, write an element of it as
$u\exp X$, with $u$ unitary and $X$ a triple of Hermitian matrices.
Finiteness of the character gives $\sum_a\operatorname{tr}X_a=0$.
Its scalar part is therefore in $\operatorname{Lie}S$ and fixes $T$.
The remaining part is special-linear.  At a critical tensor, the
function $t\mapsto\|\exp(tX)T\|^2$ has a minimum at zero.  Its value
at one equals its value at zero, since $u\exp X$ fixes $T$.
Convexity, or the spectral expansion of this function, implies
$X\cdot T=0$.  Hence both factors $u$ and $\exp X$ fix $T$, and
so does the adjoint.  This is also the usual polar-decomposition
description of a critical stabilizer.  In particular
$\widetilde H^\circ$ is reductive.  The group $J$ is adjoint-closed
as well, because $h$ is unitary.

Each of the five vector groups in (K-035d) is nonzero.  Its vector
spans a line invariant under the action of $J$ on its selected
two-dimensional factor; the two outside coordinates contribute
only scalar characters.  Each factor has at least one such
invariant line.  Adjoint closure makes its orthogonal complement
invariant too.  Consequently a block-unitary change of basis,
which preserves $U$, $h$, and criticality, makes every element
of $J$ diagonal.  Thus $J$ is exactly the full diagonal stabilizer
of the resulting tensor support.

The centralizer of the effective class $\bar h$ in $H$ is exactly
$J/S$.  Indeed, a lift of an element of that centralizer satisfies
$x h x^{-1}=s h$ for some $s\in S$.  In each factor the scalar
$s_a$ preserves the spectrum of $h$.  Its unique double eigenvalue
must be fixed, so $s_a=1$.  Thus $x$ centralizes the actual $h$.

\begin{lemma}[K-035b]
If $H$ has positive dimension, then $J/S$ has positive dimension.
\end{lemma}

\begin{proof}
Conjugation by $h$ is a finite-order automorphism of the connected
reductive group $\widetilde H^\circ$.  It preserves a maximal torus
$Q$ of that group, by the classical stable Borel--torus theorem;
the precise source used below is Serre, Section 3.3, Proposition
3(c).  Its action on the $Q$-weights in each natural factor has
orbits of two-power length.  A nontrivial orbit has even length.
On the sum of its weight spaces, $h$ is a cyclic block operator,
whose eigenvalue multiset is invariant under multiplication by
$-1$: its characteristic polynomial is a polynomial in an even
power of the variable.  Such an orbit would give an antipodal
pair in that factor's $h$-spectrum.

The first and third spectra in (K-035b) are antipodal-free.  Hence
every $Q$-weight occurring in the first or third natural factor
is fixed by $h$.  Their combined representation is faithful on
$\widetilde H$.  Indeed, an element acting identically on the
first and third factors and fixing $T$ must act identically on
the second factor, because the tensor is concise; identity
reduced matrices imply that conciseness.  These weights therefore
span $X^*(Q)\otimes\mathbb Q$.  The automorphism of $Q$ induced
by $h$ is the identity.  Thus $Q\subseteq J$.

The central torus $S$ has dimension two and is contained in $Q$.
If $\dim H>0$, the connected reductive quotient
$\widetilde H^\circ/S$ is positive-dimensional and has positive
rank.  Therefore $\dim Q>2$, and $\dim J/S>0$ as asserted.
\end{proof}

\subsection*{The case of a finite full effective stabilizer}

Suppose first that $H$ is finite.  After the block-unitary
diagonalization just established, the actual support has incidence
rank $3\cdot10-2=28$, because its full diagonal stabilizer is $J$
and $J/S$ is finite.  Choose one supported coordinate in each of
the $25$ old edge groups as a base.  Their columns are independent:
after collapsing the selected coordinates they are the independent
old columns.  Extend them to $28$ independent supported columns.

Apply the integral row transformation which replaces the old
selected row in each factor by the sum of the old and new rows,
and retains the new row as an extra row.  Delete the old last
row in the second and third factors.  The old $25\times25$
incidence minor has absolute determinant $512$.  Subtract each
base column from the additional columns in its group.  The new
$28\times28$ minor is block triangular, of determinant
\[
                            \pm512\det D,               \tag{K-035e}
\]
where $D$ has three columns.  Core-group columns are differences
of binary cube vertices from a common base vertex, and a vector
group contributes a signed coordinate unit.  Row signs determined
by the core base, and independent column signs on unit columns,
turn $D$ into a binary $3\times3$ matrix.  The elementary
Hadamard argument of K-034 gives $|\det D|\le2$.

The order of the finite effective diagonal stabilizer is the index
of the support lattice in its saturation, and therefore divides
every nonzero maximal incidence minor.  The two block-sum relations
are primitive, so the row deletion above computes a minor in an
integral basis of the saturated incidence space.  We conclude that
\[
                     |C_H(\bar h)|=|J/S|\le1024.        \tag{K-035f}
\]

For any finite group with character $\chi$, the conjugacy class
of an element lies in one character coset.  Applying this to
$\bar h$ gives
\[
 \frac{|H|}{|C_H(\bar h)|}\le |\ker\chi|,
 \qquad\hbox{hence}\qquad
                         |\chi(H)|\le |C_H(\bar h)|.
                                                               \tag{K-035g}
\]
But $\chi(h)$ has order $512$.  The finite subgroup
$\chi(H)\subset\mathbb C^*$ consequently has order divisible by
$512$ and at most $1024$.  It is $\mu_{512}$ or $\mu_{1024}$.

\subsection*{An exhaustive small balance calculation}

We now prepare the case $\dim J/S>0$.  After its block-unitary
diagonalization, choose a non-scalar cocharacter of $J$.
Subtract a cocharacter of $S$ so that its weights at the first
coordinate of the first and second factors are zero.  This fixes
the two scalar freedoms.  Write its remaining weight vector as $x$.

The squared moduli of the tensor coordinates give a nonnegative
balance vector.  The scalar coordinates have their fixed positive
masses.  The six nonscalar groups contain $8+2+2+2+2+2=18$
variables, with totals (K-035d).  Their three additional equations
say that the mass of the new coordinate in each selected block
is one.  Choose a vertex of the face obtained by setting all
absent tensor coordinates to zero.  It is a vertex of the full
balance polytope, and its support is contained in that of $T$.
Therefore $x$ annihilates every incidence column in this vertex.

Here is the exact finite classification used in the proof:
\[
\begin{array}{c|r|c}
\text{positive nonscalar coordinates}&\text{vertices}
 &\text{kernel dimension modulo }S\\ \hline
9&472&0\\
8&56&1\\
7&8&2
\end{array}                                                   \tag{K-035h}
\]
In the middle row the kernel is spanned by
$E_{a,i}-E_{b,j}$, with $a\ne b$.  In the last row it is
\[
 \{u_AE_{A,i}+u_BE_{B,j}+u_CE_{C,k}:u_A+u_B+u_C=0\}.
\]
Here every indicated coordinate belongs to its selected
two-dimensional block, and $E_{a,i}$ gives weight one on that
coordinate and zero elsewhere.  In the last row, the core has
exactly two nonzero entries, at opposite cube vertices, of masses
$293/512$ and $1$; every vector is on the side of the $293/512$
entry.  The isolated coordinates in the table belong to the
mass-one entry.  In the middle row the three possible pairs of
factors occur $24,8,24$ times for $AB,AC,BC$, respectively.

For reproducibility, the full exhaustive procedure is as follows.
Give the cube its eight binary feature vectors, and each vector
group the features $0,e_a$.  There are $8\cdot2^5=256$ ways to
choose one base variable in each group.  Eliminate those six
variables with the group totals.  Among the remaining twelve
columns choose three, giving $256\binom{12}{3}=56\,320$ choices.
For each nonzero determinant, solve the three remaining equations
by integer Cramer's rule and test all eighteen nonnegativity
conditions.  The three-row determinant is always $\pm1$ or $\pm2$.
Thus integer numerators over $1024$ suffice throughout.

This enumerates every vertex.  Every group has positive total, so
a vertex has a positive base in each group.  After elimination,
its other positive columns are independent and number at most
three.  They extend to a basis of the three-row column space,
since the vector groups supply the three coordinate directions.
The added basis variables have value zero.  Hence the vertex
occurs in the procedure just specified.

There are exactly $32\,000$ nonsingular choices, of which $6080$
are feasible with repetitions; they give $536$ distinct vertices.
The retained verifier independently checks the rank of every
positive reduced-column set.  For the $64$ degenerate vertices
it computes the kernel of the full support-incidence transpose,
with the two scalar gauges imposed, over the rational numbers.
These kernels give exactly (K-035h).  All $536$ mass vectors and
all $64$ kernel certificates are retained, not only the counts:
\texttt{verify\_rank\_two\_boundary\_stabilizers.py} and
\texttt{rank\_two\_boundary\_stabilizers\_checks.json}.
The script also checks the old minor, all $1000$ fixed weights,
the group allocations, and the antipodal assertions in (K-035b).
This is a finite exact calculation, with no floating-point
feasibility test or random search.

\subsection*{Criticality forces a direct sum}

The first row of (K-035h) is impossible for the nonzero gauged
cocharacter $x$.  In its second row, after permuting factors for
notation, $x$ gives opposite nonzero weights to lines
$B_1\subset B_2$ and $C_1\subset C_2$, and weight zero elsewhere.
Its invariant tensor must have the form
\[
       T=T_0+a\otimes b_1\otimes c_1,
 \qquad T_0\in A\otimes B_1^\perp\otimes C_1^\perp,
                                                               \tag{K-035i}
\]
with $b_1,c_1$ unit vectors.  The complete fixed support implies
$a\in A_2$: only the core group uses both selected second and
third coordinates.  The identity reduced matrices in the second
and third factors give $\|a\|=1$.  The first reduced matrix is
\[
                         R_A(T_0)+aa^*=I_A.
\]
Since $R_A(T_0)$ is positive semidefinite, its value on the unit
vector $a$ is zero.  Contraction of $T_0$ with $a^*$ is therefore
zero.  Hence $T_0\in a^\perp\otimes B_1^\perp\otimes C_1^\perp$.
The tensor is an orthogonal direct sum of a concise
$9\times9\times9$ tensor and a nonzero one-dimensional tensor.

In the last row of (K-035h), the cocharacter gives weights
$u_A,u_B,u_C$ to three selected lines, with sum zero, and is
zero elsewhere.  If one weight is zero, this is the preceding
case.  Otherwise no nonempty proper subset of the three weights
has sum zero.  Invariance immediately gives the same direct-sum
decomposition, with all three indicated lines isolated.

In either case, the remaining selected line in each
two-dimensional block is the orthogonal complement of the
isolated line.  Every frame vector lies in that remaining line,
and all nineteen scalar coordinates remain nonzero.  The
one-dimensional summand has squared norm one.  Subtracting its
core mass from (K-035c) leaves exactly the original positive
mass vector $p$ on the old support $\Sigma$.  Thus the
nine-dimensional summand has precisely one nonzero coefficient
on each edge of $\Sigma$ and no other coordinates.  Independence
of the old incidence columns makes the corresponding diagonal
torus map onto the torus of these $25$ nonzero coefficients.
All such tensors are diagonally equivalent.  We have proved
\[
               \dim J/S>0\quad\Longrightarrow\quad
                         T\simeq T_9\oplus T_1,         \tag{K-035j}
\]
where $T_1$ is a nonzero one-dimensional tensor.

\subsection*{The full stabilizer of the direct sum}

We verify that no full stabilizer can mix the two summands in
(K-035j).  For a concise tensor $R$, consider the algebra of
triples of endomorphisms satisfying
\[
 (a\otimes1\otimes1)R=(1\otimes b\otimes1)R
                         =(1\otimes1\otimes c)R.        \tag{K-035k}
\]
For $T_9$ this algebra is the scalar diagonal copy of $\mathbb C$:
the equality of the first two actions gives an infinitesimal
stabilizer $(a,-b,0)$, and the equality with the third gives
$(a,0,-c)$.  Its previously proved Lie stabilizer $S_9$ forces
$a=b=c$ scalar.

For $T_9\oplus T_1$, the off-diagonal blocks in (K-035k) vanish.
For example, an off-diagonal first-factor block produces a tensor
in $A_1\otimes B_9\otimes C_9$ or
$A_9\otimes B_1\otimes C_1$.  Neither component can be produced
by acting in the second factor.  Conciseness of the two summands
therefore makes each such block zero.  The other factors are
the same.  The algebra in (K-035k) is consequently
$\mathbb C\oplus\mathbb C$, with its two primitive idempotents
the simultaneous projections onto the two tensor summands.

The full stabilizer acts on this algebra by factorwise
conjugation.  It cannot exchange the idempotents, since their
ranks in each factor are $9$ and $1$.  It preserves both
summands, and its determinant character is the product of the
characters of their stabilizers.  The one-dimensional tensor
has determinant-product character one on its stabilizer, whereas
K-020 gives character image $\mu_{512}$ on $T_9$.  Hence
\[
                 \chi(\operatorname{Stab}(T_9\oplus T_1))
                               =\mu_{512}.              \tag{K-035l}
\]

If $H$ has positive dimension, Lemma K-035b and (K-035j)--(K-035l)
apply.  If $H$ is finite, (K-035f)--(K-035g) apply.  In both
cases its character image is $\mu_{512}$ or $\mu_{1024}$.

Finally this character image is exactly the grading stabilizer
of the quotient point.  If a grading element fixes that point,
choose a full product-linear element with that character.  Its
translate of the closed special-linear orbit is the same closed
orbit, so an element of $SL_{10}^3$ corrects it to an actual
stabilizer element.  The converse is immediate.  This proves
(K-035a) at every nonzero point of $Y$.

\subsection*{Scope in the global cancellation problem}

K-035a excludes every higher odd grading stabilizer on $Y$,
including branch-permutation inertia over its conductor.  The
finite-length Koszul argument and cyclotomic no-collision lemma
of K-027 now apply directly to every graded coherent module on
$Y$, and likewise on its cyclic translates.  Their transported
sum is three times a three-integral Laurent series.  Dividing by
the moving-generator factor of K-027 preserves this conclusion,
because that factor is a three-adic unit at the primitive
dyadic root.  Thus none of these three component closures occurs
among the components indexed in the union defining
$Z_{\mathrm{coll}}$ in K-031.  Intersections with other collision
component closures are not asserted to be empty.

The conductor itself remains uncomputed, and no isomorphism with
the normalization is asserted.  The other collision components,
including the actual large-rank examples of K-029, remain
uncontrolled.  In particular this result does not prove an
actual global pole at order $2^{n-1}$ for arbitrary $n$, and
does not settle a strict positivity or saturation index.

\paragraph{Primary source and verification scope.}
The stable maximal torus used in Lemma K-035b is classical:
Serre, J.-P. (2007). Bounds for the orders of the finite subgroups
of $G(k)$. In M. Geck, D. Testerman, and J. Th\'evenaz (Eds.),
\emph{Group representation theory} (pp. 405--450). EPFL Press.
The author-hosted manuscript, Section 3.3, Proposition 3(c),
printed pp. 13--14, has been read together with the complete
proposition and its proof sketch.  It credits Springer--Steinberg
(1968), I.2.9, for that assertion; no reading of that original
proof is claimed here.  The finite-group character bound
(K-035g) is the elementary conjugacy-coset argument displayed
above and uses no character-table classification.  All new
finite arithmetic is retained in the exact verifier and full
certificate named above.  The geometric proof treats the natural
factor representations only on the full stabilizer, and separately
proves the effective centralizer identity before using its order.
